\scanpage{210}
\mspage{84}
dans $\underline{\mathrm{Fl}}(M)$. Comme $\underrightarrow{\lim}\gamma_i$ \deleted{$\unsure{L|i\in J}$} \deleted{$\unsure{=B}$} est un iso, il est dans $C$. \deleted{Car la stabilité de $C$ implique que $C$} Comme $C$ est $\unsure{\pi|R}$-accessible, il s'ensuit que l'ens $J'$ formé des \deleted{$\unsure{J}$} $i\in J$ tels que $\gamma_i\in C$ est lui aussi cofinal.$^{+}$
% slanted
\margin{\uuline{canulé}, cf annotation précédente.} Donc on trouve $i$ avec $\alpha_i\in C^\pi$, $\gamma_i\in C$, ce qui établit le lemme 6.4, et achève la démonstration du théorème 6.1,

\deleted{Chemin faisant j'ai de plus en plus l'impression que le rôle que je fais jouer dans \unsure{tout ça|tout §ci} aux monomorphismes, et surtout l'hypothèse que $C$ soit formé de monomorphismes, est tout à fait artificiel. Je vais essayer de m'en débarrasser.}

7. \uline{Construction de triples de Quillen clos avec} $C=\mathrm{Mon}\subset\mathrm{Cof}_W$

Finalement, tous les beaux développements des §§ 5,6 ne me donnent pas le ``sweeping result'' souhaité. Peut-être est-il faux --- peut-être la propriété d'engendrement $\sim$ par une petite partie (pour les ensembles accessibles, ou L-accessibles) est-elle beaucoup plus \unsure{rarement|souvent} satisfaite, que je ne l'aurais imaginé. Pour bien faire, il faudrait que je traite jusqu'au
% archive sheet number: 209
